\cdpchapter{Resumen}

Este Trabajo Fin de Grado presenta el dise\~no, implementaci\'on y validaci\'on
de un sistema web para la gesti\'on de entregables acad\'emicos,
orientado a mejorar la operativa diaria de profesorado y alumnado.
La soluci\'on cubre el ciclo completo de trabajo:
gesti\'on de cursos y actividades, configuraci\'on de entregables,
realizaci\'on de entregas por parte del estudiante,
evaluaci\'on docente con feedback y control de visibilidad de notas.

Desde el punto de vista t\'ecnico, la plataforma se implementa con
arquitectura cliente-servidor desacoplada,
empleando frontend en React + TypeScript y backend en Spring Boot sobre Java 21,
con persistencia relacional y migraciones versionadas mediante Flyway.
La seguridad se basa en autenticaci\'on JWT con renovaci\'on de sesi\'on,
autorizaci\'on por rol y validaciones por contexto acad\'emico.

El proyecto integra adem\'as una estrategia de calidad en capas
(pruebas unitarias, integraci\'on, E2E y mutaci\'on),
junto con automatizaci\'on CI/CD y despliegue en entorno cloud,
lo que permite reducir regresiones y mejorar trazabilidad de cambios.

Como resultado, se obtiene una base funcional y mantenible,
alineada con los requisitos identificados en el an\'alisis inicial
y con capacidad real de evoluci\'on hacia funcionalidades avanzadas
de seguridad, observabilidad e interoperabilidad.

\cdpchapter{Abstract}

This Bachelor's Thesis presents the design, implementation, and validation
of a web-based academic assignment management platform
aimed at improving the day-to-day workflow of both lecturers and students.
The system supports the full process:
course and activity management, assignment configuration,
student submissions, lecturer grading with feedback,
and controlled grade visibility.

From a technical perspective, the platform follows a decoupled
client-server architecture,
using a React + TypeScript frontend and a Java 21 Spring Boot backend,
with relational persistence and versioned schema migrations via Flyway.
Security is implemented through JWT-based authentication,
session refresh, role-based authorization, and context-aware access checks.

The project also incorporates a layered quality strategy
(unit, integration, end-to-end, and mutation testing),
together with CI/CD automation and cloud deployment,
thus reducing regression risk and improving traceability.

The outcome is a functional and maintainable solution,
aligned with the initial requirements and ready to evolve
towards more advanced capabilities in security, observability,
and institutional interoperability.


